\section{Introdução}

\paragraph{}
Este documento apresenta as evidências práticas das atividades desenvolvidas na
Aula 11 do módulo de Introdução ao Pentest, no âmbito do Programa Hackers do
Bem -- Nível Fundamental. A aula teve como objetivo aprofundar o entendimento
sobre mecanismos internos de sistemas Linux relacionados ao gerenciamento de
usuários, permissões de arquivos, execução privilegiada de programas, tarefas
agendadas e técnicas de reconhecimento local, conhecimentos fundamentais para a
compreensão de processos de escalação de privilégios em ambientes Linux.

\paragraph{}
Conforme as orientações propostas, este material reúne os registros das
atividades realizadas durante a aula, com destaque para a análise do mecanismo
SUID (\textit{Set User ID}), incluindo a criação, compilação e execução de
programas com privilégios diferenciados. Além disso, foram explorados conceitos
relacionados ao gerenciamento de contas de usuário por meio dos arquivos
\texttt{/etc/passwd} e \texttt{/etc/shadow}, ao funcionamento de tarefas
agendadas utilizando \texttt{cron} e \texttt{crontab}, bem como à coleta de
informações do sistema por meio de técnicas de reconhecimento local e da
ferramenta \texttt{pspy}. Como evidência da atividade, são apresentados os
resultados dos comandos solicitados no roteiro: visualização do código-fonte
\texttt{suid.c}, listagem das permissões dos arquivos relacionados e execução
do binário \texttt{suid}.
